\documentclass[11pt,letterpaper]{article}
\usepackage[technical-reference]{cornell-notes}
\usepackage{needspace}

\title{Clause 10: Modules - Cornell Notes}
\author{}
\date{}
\setDocTitle{Clause 10: Modules}
\setDocSubtitle{C++ 2024 Cornell Notes}
\setDocOwner{C++ 2024 Cornell Notes}
\setDocAuthor{}
\setDocDate{}
\setCornellCollection{C++ 2024 Cornell Notes}
\setCornellUnitType{Clause}
\setCornellUnitNumber{10}
\setCornellUnitTitle{Modules}
\setCornellCompanionLabel{C++ Technical Reference}
\hypersetup{pdftitle={Clause 10: Modules - Cornell Notes},pdfsubject={C++ technical study notes},pdfauthor={}}

\begin{document}
\maketitle
\begin{CornellReferenceOrientationBox}
\textbf{Purpose.} Defines module units and purviews, exports, imports, global and private fragments, instantiation context, visibility, and reachability.
\par\textbf{Role.} Provides a semantic distribution and visibility model beyond textual inclusion.
\par\textbf{Principal dependencies.} Uses declarations, linkage, translation units, templates, and preprocessing boundaries.
\par\textbf{Primary audience.} programmers, build-system authors, and compiler implementers.
\par\textbf{Major subject groups.} Module units and purviews; Export declaration; Import declaration; Global module fragment; Private module fragment; Instantiation context; Reachability.
\end{CornellReferenceOrientationBox}
\section{Learning Objectives}
\begin{itemize}[label=\textcolor{CornellReferenceTeal}{$\blacktriangleright$}]
\item Define the governing ideas of Module units and purviews and state the consequences for a program or implementation.
\item Identify the governing ideas of Export declaration and state the consequences for a program or implementation.
\item Distinguish the governing ideas of Import declaration and state the consequences for a program or implementation.
\item Explain the governing ideas of Global module fragment and state the consequences for a program or implementation.
\item Trace the governing ideas of Private module fragment and state the consequences for a program or implementation.
\end{itemize}
\section{Normative Structure Map}
\small\begin{longtable}{@{}>{\RaggedRight\arraybackslash}p{0.13\textwidth}>{\RaggedRight\arraybackslash}p{0.27\textwidth}>{\RaggedRight\arraybackslash}p{0.19\textwidth}>{\RaggedRight\arraybackslash}p{0.31\textwidth}@{}}
\toprule\textbf{Subclause} & \textbf{Subject} & \textbf{Rule category} & \textbf{Key dependencies / entities}\\\midrule\endfirsthead
\toprule\textbf{Subclause} & \textbf{Subject} & \textbf{Rule category} & \textbf{Key dependencies / entities}\\\midrule\endhead
\bottomrule\endfoot
10.1 & Module units and purviews & core-language semantics & Uses declarations, linkage, translation units, templates, and preprocessing boundaries.\\
10.2 & Export declaration & core-language semantics & Uses declarations, linkage, translation units, templates, and preprocessing boundaries.\\
10.3 & Import declaration & core-language semantics & Uses declarations, linkage, translation units, templates, and preprocessing boundaries.\\
10.4 & Global module fragment & core-language semantics & Uses declarations, linkage, translation units, templates, and preprocessing boundaries.\\
10.5 & Private module fragment & core-language semantics & Uses declarations, linkage, translation units, templates, and preprocessing boundaries.\\
10.6 & Instantiation context & core-language semantics & Uses declarations, linkage, translation units, templates, and preprocessing boundaries.\\
10.7 & Reachability & core-language semantics & Uses declarations, linkage, translation units, templates, and preprocessing boundaries.\\
\end{longtable}\normalsize
\begin{CornellReferenceNormativeBox}A heading maps the clause; it does not replace the detailed conditions. For each operation, read requirements, effects, result, exceptions, complexity, remarks, and notes with their stated normative force.\end{CornellReferenceNormativeBox}
\subsection*{Rule Classification Guide}
\small\begin{longtable}{@{}>{\RaggedRight\arraybackslash}p{0.25\textwidth}>{\RaggedRight\arraybackslash}p{0.67\textwidth}@{}}\toprule\textbf{Label or category} & \textbf{How to read it}\\\midrule\endfirsthead\toprule\textbf{Label or category} & \textbf{How to read it (continued)}\\\midrule\endhead\bottomrule\endfoot
Well-formed / ill-formed & Formation is governed by syntax and semantic rules; an ill-formed program does not become well-formed because one implementation accepts it.\\
Diagnosable rule & A violation requires at least one diagnostic unless the wording places the case outside the diagnosable set.\\
No diagnostic required & The program can be ill-formed while the implementation has no diagnostic obligation.\\
Undefined behavior & No execution requirements are imposed; translation success does not establish safety or portability.\\
Unspecified behavior & One of the permitted outcomes may be chosen without documentation.\\
Implementation-defined behavior & The implementation chooses and documents the behavior.\\
Conditionally supported & The implementation may omit support but documents unsupported constructs.\\
\end{longtable}\normalsize
\section{Cornell Cue-and-Notes}
\Needspace{8\baselineskip}
\subsection*{10.1 Module units and purviews}
\begin{longtable}{@{}>{\RaggedRight\arraybackslash\columncolor{CornellReferenceCueGray}}p{0.28\textwidth}>{\RaggedRight\arraybackslash}p{0.64\textwidth}@{}}
\rowcolor{CornellReferenceNavy}\color{white}\textbf{Cue / recall prompt} & \color{white}\textbf{Notes}\\\endfirsthead
\rowcolor{CornellReferenceNavy}\color{white}\textbf{Cue / recall prompt} & \color{white}\textbf{Notes (continued)}\\\endhead
\arrayrulecolor{CornellReferenceLineGray}
\textbf{What rule applies to Module units and purviews?}\\[-1mm]{\scriptsize\CornellClauseRef{10.1}} & Explains the traversal or view contract for Module units and purviews; prove domain validity, lifetime, sentinel compatibility, and any borrowing or invalidation property.\\\hline
\end{longtable}
\begin{CornellReferenceSummaryBox}Module units and purviews supplies the core-language semantics portion of this clause. The key review move is to connect its local wording to the clause-wide model, then classify each condition by normative force before relying on it.\end{CornellReferenceSummaryBox}
\Needspace{8\baselineskip}
\subsection*{10.2 Export declaration}
\begin{longtable}{@{}>{\RaggedRight\arraybackslash\columncolor{CornellReferenceCueGray}}p{0.28\textwidth}>{\RaggedRight\arraybackslash}p{0.64\textwidth}@{}}
\rowcolor{CornellReferenceNavy}\color{white}\textbf{Cue / recall prompt} & \color{white}\textbf{Notes}\\\endfirsthead
\rowcolor{CornellReferenceNavy}\color{white}\textbf{Cue / recall prompt} & \color{white}\textbf{Notes (continued)}\\\endhead
\arrayrulecolor{CornellReferenceLineGray}
\textbf{What rule applies to Export declaration?}\\[-1mm]{\scriptsize\CornellClauseRef{10.2}} & Defines the core-language rules for Export declaration, including formation, semantic effect, interactions with type and lookup, and any ill-formed or implementation-dependent cases.\\\hline
\end{longtable}
\begin{CornellReferenceSummaryBox}Export declaration supplies the core-language semantics portion of this clause. The key review move is to connect its local wording to the clause-wide model, then classify each condition by normative force before relying on it.\end{CornellReferenceSummaryBox}
\Needspace{8\baselineskip}
\subsection*{10.3 Import declaration}
\begin{longtable}{@{}>{\RaggedRight\arraybackslash\columncolor{CornellReferenceCueGray}}p{0.28\textwidth}>{\RaggedRight\arraybackslash}p{0.64\textwidth}@{}}
\rowcolor{CornellReferenceNavy}\color{white}\textbf{Cue / recall prompt} & \color{white}\textbf{Notes}\\\endfirsthead
\rowcolor{CornellReferenceNavy}\color{white}\textbf{Cue / recall prompt} & \color{white}\textbf{Notes (continued)}\\\endhead
\arrayrulecolor{CornellReferenceLineGray}
\textbf{What rule applies to Import declaration?}\\[-1mm]{\scriptsize\CornellClauseRef{10.3}} & Defines the core-language rules for Import declaration, including formation, semantic effect, interactions with type and lookup, and any ill-formed or implementation-dependent cases.\\\hline
\end{longtable}
\begin{CornellReferenceSummaryBox}Import declaration supplies the core-language semantics portion of this clause. The key review move is to connect its local wording to the clause-wide model, then classify each condition by normative force before relying on it.\end{CornellReferenceSummaryBox}
\Needspace{8\baselineskip}
\subsection*{10.4 Global module fragment}
\begin{longtable}{@{}>{\RaggedRight\arraybackslash\columncolor{CornellReferenceCueGray}}p{0.28\textwidth}>{\RaggedRight\arraybackslash}p{0.64\textwidth}@{}}
\rowcolor{CornellReferenceNavy}\color{white}\textbf{Cue / recall prompt} & \color{white}\textbf{Notes}\\\endfirsthead
\rowcolor{CornellReferenceNavy}\color{white}\textbf{Cue / recall prompt} & \color{white}\textbf{Notes (continued)}\\\endhead
\arrayrulecolor{CornellReferenceLineGray}
\textbf{What rule applies to Global module fragment?}\\[-1mm]{\scriptsize\CornellClauseRef{10.4}} & Places Global module fragment within the module ownership and dependency model; visibility, reachability, linkage, and textual preprocessing remain distinct.\\\hline
\end{longtable}
\begin{CornellReferenceSummaryBox}Global module fragment supplies the core-language semantics portion of this clause. The key review move is to connect its local wording to the clause-wide model, then classify each condition by normative force before relying on it.\end{CornellReferenceSummaryBox}
\Needspace{8\baselineskip}
\subsection*{10.5 Private module fragment}
\begin{longtable}{@{}>{\RaggedRight\arraybackslash\columncolor{CornellReferenceCueGray}}p{0.28\textwidth}>{\RaggedRight\arraybackslash}p{0.64\textwidth}@{}}
\rowcolor{CornellReferenceNavy}\color{white}\textbf{Cue / recall prompt} & \color{white}\textbf{Notes}\\\endfirsthead
\rowcolor{CornellReferenceNavy}\color{white}\textbf{Cue / recall prompt} & \color{white}\textbf{Notes (continued)}\\\endhead
\arrayrulecolor{CornellReferenceLineGray}
\textbf{What rule applies to Private module fragment?}\\[-1mm]{\scriptsize\CornellClauseRef{10.5}} & Places Private module fragment within the module ownership and dependency model; visibility, reachability, linkage, and textual preprocessing remain distinct.\\\hline
\end{longtable}
\begin{CornellReferenceSummaryBox}Private module fragment supplies the core-language semantics portion of this clause. The key review move is to connect its local wording to the clause-wide model, then classify each condition by normative force before relying on it.\end{CornellReferenceSummaryBox}
\Needspace{8\baselineskip}
\subsection*{10.6 Instantiation context}
\begin{longtable}{@{}>{\RaggedRight\arraybackslash\columncolor{CornellReferenceCueGray}}p{0.28\textwidth}>{\RaggedRight\arraybackslash}p{0.64\textwidth}@{}}
\rowcolor{CornellReferenceNavy}\color{white}\textbf{Cue / recall prompt} & \color{white}\textbf{Notes}\\\endfirsthead
\rowcolor{CornellReferenceNavy}\color{white}\textbf{Cue / recall prompt} & \color{white}\textbf{Notes (continued)}\\\endhead
\arrayrulecolor{CornellReferenceLineGray}
\textbf{What rule applies to Instantiation context?}\\[-1mm]{\scriptsize\CornellClauseRef{10.6}} & Defines the core-language rules for Instantiation context, including formation, semantic effect, interactions with type and lookup, and any ill-formed or implementation-dependent cases.\\\hline
\end{longtable}
\begin{CornellReferenceSummaryBox}Instantiation context supplies the core-language semantics portion of this clause. The key review move is to connect its local wording to the clause-wide model, then classify each condition by normative force before relying on it.\end{CornellReferenceSummaryBox}
\Needspace{8\baselineskip}
\subsection*{10.7 Reachability}
\begin{longtable}{@{}>{\RaggedRight\arraybackslash\columncolor{CornellReferenceCueGray}}p{0.28\textwidth}>{\RaggedRight\arraybackslash}p{0.64\textwidth}@{}}
\rowcolor{CornellReferenceNavy}\color{white}\textbf{Cue / recall prompt} & \color{white}\textbf{Notes}\\\endfirsthead
\rowcolor{CornellReferenceNavy}\color{white}\textbf{Cue / recall prompt} & \color{white}\textbf{Notes (continued)}\\\endhead
\arrayrulecolor{CornellReferenceLineGray}
\textbf{What rule applies to Reachability?}\\[-1mm]{\scriptsize\CornellClauseRef{10.7}} & Places Reachability within the module ownership and dependency model; visibility, reachability, linkage, and textual preprocessing remain distinct.\\\hline
\end{longtable}
\begin{CornellReferenceSummaryBox}Reachability supplies the core-language semantics portion of this clause. The key review move is to connect its local wording to the clause-wide model, then classify each condition by normative force before relying on it.\end{CornellReferenceSummaryBox}
\section{Language-Lawyer and Library-Usage Pitfalls}
\begin{CornellReferencePitfallBox}\begin{itemize}[label=\textcolor{CornellReferenceAmber}{$\blacktriangleright$}]
\item Treating import as textual inclusion.
\item Confusing visibility with reachability.
\item Exporting a declaration while overlooking non-reachable dependencies in its interface.
\item Placing preprocessing-dependent material in the wrong fragment.
\item Treating a note, example, or recommended practice as though it imposed a mandatory program rule.
\end{itemize}\end{CornellReferencePitfallBox}
\section{Programmer and Implementation Checklist}
\begin{CornellReferenceChecklistBox}\begin{itemize}[label=$\square$]
\item Classify every unit as interface, implementation, partition, or non-module unit.
\item Mark the purview and fragment boundaries.
\item Trace exported declarations and their semantic dependencies.
\item Check import ordering and permitted forms.
\item Validate instantiation contexts and reachable definitions.
\item For \CornellClauseRef{10.1}, verify module units and purviews against the precise rule category and all cited dependencies.
\item For \CornellClauseRef{10.2}, verify export declaration against the precise rule category and all cited dependencies.
\item For \CornellClauseRef{10.3}, verify import declaration against the precise rule category and all cited dependencies.
\item For \CornellClauseRef{10.4}, verify global module fragment against the precise rule category and all cited dependencies.
\end{itemize}\end{CornellReferenceChecklistBox}
\section{Clause Synthesis}
\begin{CornellReferenceOrientationBox}Defines module units and purviews, exports, imports, global and private fragments, instantiation context, visibility, and reachability. Provides a semantic distribution and visibility model beyond textual inclusion. A sound review distinguishes syntax and participation from runtime preconditions, keeps notes and examples non-mandatory, and follows cross-clause dependencies before making a portability claim.\end{CornellReferenceOrientationBox}
\section{Self-Test Questions}
\begin{enumerate}
\item What role does Module units and purviews play in the clause's overall model?
\item Which normative distinctions must be kept separate when applying Export declaration?
\item How would you audit a program or implementation that depends on Import declaration?
\item Scenario: a program relies on an unstated implementation behavior while using Global module fragment. What must be checked before calling the program portable?
\item Scenario: a program relies on an unstated implementation behavior while using Private module fragment. What must be checked before calling the program portable?
\end{enumerate}
\clearpage\section{Answer Key}
\begin{enumerate}
\item Module units and purviews is one part of the clause's core-language semantics model. Its role is understood by connecting the local rule in 10.1 to the clause orientation and its dependent concepts.
\item Keep formation and well-formedness separate from runtime preconditions and effects. Also distinguish mandatory wording from notes, implementation choice from undefined behavior, and interface declarations from detailed contracts; see 10.2.
\item Audit Import declaration by locating the controlling wording in 10.3, listing inputs and type constraints, proving any preconditions, recording effects and failure behavior, and checking lifetime, invalidation, complexity, and synchronization where applicable.
\item First identify the exact requirement in 10.4, then classify the relied-on behavior as required, unspecified, implementation-defined, conditionally supported, or undefined. Portability requires satisfying the program obligations and avoiding dependence on a choice not guaranteed in the target environments.
\item First identify the exact requirement in 10.5, then classify the relied-on behavior as required, unspecified, implementation-defined, conditionally supported, or undefined. Portability requires satisfying the program obligations and avoiding dependence on a choice not guaranteed in the target environments.
\end{enumerate}
\section{Key-Term Recap}
\begin{longtable}{@{}>{\RaggedRight\arraybackslash}p{0.28\textwidth}>{\RaggedRight\arraybackslash}p{0.64\textwidth}@{}}\toprule\textbf{Term} & \textbf{Working meaning}\\\midrule\endfirsthead\toprule\textbf{Term} & \textbf{Working meaning (continued)}\\\midrule\endhead\bottomrule\endfoot
module unit & A translation unit associated with a named module or module partition.\\
module purview & The portion of a unit governed by its module declaration.\\
export & A declaration that makes selected entities reachable to importers under the module rules.\\
import & A declaration establishing visibility and semantic dependency on another unit.\\
global module fragment & A leading region for material before the named module purview.\\
private module fragment & A region confined to the current module unit.\\
reachability & Whether a declaration or definition can contribute to semantic interpretation in a context.\\
\end{longtable}
\CornellTakeaway{Modules control semantic ownership and reachability, so interfaces must expose every dependency needed to interpret exported declarations.}
\end{document}
